% ============================================================
% Отчёт по исследованию ЯМ — баскетбол
% Адиатуллин Т. Р., СПбПУ, 2026
% ============================================================

\documentclass[a4paper, final]{article}

\usepackage[russian]{babel}
\usepackage{fontspec}
\setmainfont{Times New Roman}
\newfontfamily\cyrillicfonttt{Courier New}
\usepackage[14pt]{extsizes}
\usepackage[left=20mm, top=20mm, right=20mm, bottom=20mm, footskip=10mm]{geometry}
\usepackage{ragged2e}
\usepackage{indentfirst}
\usepackage{graphicx}
\usepackage{tikz}
\usepackage{caption}
\usepackage{xcolor}
\usepackage{subcaption}
\usepackage{booktabs}
\usepackage{hyperref}
\usepackage{array}
\usepackage{tabularx}
\usepackage{longtable}
\usepackage{multirow}
\usepackage{pdflscape}
\usepackage{soul}
\usepackage{enumitem}

\definecolor{link_color}{HTML}{0000EE}

\hypersetup{
    colorlinks=true,
    linkcolor=link_color,
    urlcolor=link_color,
    citecolor=link_color
}

% ============================================================
% Вспомогательные команды
% ============================================================

\newcommand{\scoretablefilled}[8]{%
    \vspace{4pt}
    \begin{table}[h!]
    \centering
    \small
    \begin{tabular}{|l|c|c|}
    \hline
    \textbf{Критерий (1--5)} & \textbf{Grok} & \textbf{Qwen} \\
    \hline
    Достоверность содержания & #1 & #2 \\
    \hline
    Полнота и структурированность & #3 & #4 \\
    \hline
    Уместность стиля & #5 & #6 \\
    \hline
    Искажения / ошибки (да/нет) & #7 & #8 \\
    \hline
    \end{tabular}
    \caption{Оценка ответов}
    \end{table}
    \vspace{2pt}
}

% ============================================================
\begin{document}

% ============================================================
% Титульная страница
% ============================================================
{
\thispagestyle{empty}
\begin{center}
    \hfill \break
    \hfill \break
    МИНИСТЕРСТВО НАУКИ И ВЫСШЕГО ОБРАЗОВАНИЯ РОССИЙСКОЙ ФЕДЕРАЦИИ\\[5pt]
    Федеральное государственное автономное образовательное учреждение высшего образования
    «Санкт-Петербургский политехнический университет Петра Великого»\\[10pt]
    Институт компьютерных наук и кибербезопасности\\[5pt]
    Направление: 02.03.01 Математика и компьютерные науки\\[30pt]

    {\large Отчет по проекту: «Исследование смысловой устойчивости больших языковых моделей»}\\[30pt]
    \LARGE\bf Предметная область: «Баскетбол».
\end{center}

\vspace{110pt}

\begin{tabular*}{460pt}{@{\extracolsep{\fill}} l r}
     Студент,\tabularnewline группы 5130201/40003 \hspace{50pt} & \line(1,0){100} \quad Адиатуллин Т. Р.\\
     & \\
     Руководитель,\tabularnewline Преподаватель & \line(1,0){121.5} \quad [ФИО руководителя]\\
     & \\
     & \\
     & \\
     & \guillemotleft \line(1,0){30} \guillemotright \line(1,0){100} \, 2026 \,г.
\end{tabular*} \\

\vfill

\begin{center}
    Санкт-Петербург, 2026
\end{center}
\newpage
}

\tableofcontents
\newpage

% ============================================================
\section{Постановка задачи}

В данной работе исследуется влияние функциональных стилей русского языка и речевых помех на качество ответов больших языковых моделей в предметной области \textbf{«Баскетбол»}.

Для исследования выбраны две модели: \textbf{Grok} (xAI) и \textbf{Qwen} (Alibaba). Исследуются два базовых вопроса:
\begin{enumerate}
    \item \textbf{Вопрос о тактике} — пик-н-ролл: механизм выполнения и роли игроков.
    \item \textbf{Вопрос о правилах} — тайм-аут: количество, продолжительность и порядок запроса.
\end{enumerate}

Каждый вопрос формулируется в девяти вариантах: пяти функциональных стилях и четырёх типах речевых помех. Ответы оцениваются по трём критериям:
\begin{enumerate}
    \item \textbf{Достоверность содержания} — субъективная экспертная оценка точности фактического содержания ответа (1–5).
    \item \textbf{Полнота и структурированность} — оценка полноты раскрытия темы и логичности изложения (1–5).
    \item \textbf{Уместность стиля} — оценка соответствия ответа функциональному стилю запроса согласно учебнику Д.\,Э.~Розенталя «Практическая стилистика русского языка» (1–5).
\end{enumerate}

Дополнительно фиксируются \textbf{искажения смысла и ошибки}, возникшие в ответе вследствие помех в запросе.

\newpage

% ============================================================
\section{Вопрос 1: Тактика — пик-н-ролл}
% ============================================================

% ------------------------------------------------------------
\subsection{Научный стиль}

\noindent\textbf{Промпт:} Дайте системное описание тактического взаимодействия «заслон-выход» (пик-н-ролл) в баскетболе: механизм выполнения, роли участников и условия эффективного применения.

\noindent\textbf{Ответ Grok:} В баскетболе тактическое взаимодействие заслон-выход или пик-н-ролл строится на взаимодействии двух игроков: один ставит заслон мячу, второй с мячом использует его для прохода или броска. Механизм начинается с того что игрок с мячом ведет дриблинг навстречу партнеру который занимает позицию сбоку от защитника и плотно ставит неподвижный заслон корпусом не двигаясь с места. Игрок с мячом проходит вплотную к заслонившему партнеру вынуждая защитника выбирать между преследованием и сменой опеки. После прохода заслоняющий разворачивается лицом к кольцу открываясь для передачи и потенциального броска или прохода под щит. Роли четко распределены: первый номер ведет атаку и решает бросок проход или передачу второй создает пространство и получает опцию на свободный бросок. Эффективность достигается при хорошей координации быстром чтении защиты точном тайминге движения и умении обоих игроков завершать атаку в случае свитча или дабл-тим.

\noindent\textbf{Ответ Qwen:} Тактическое взаимодействие заслон-выход представляет собой фундаментальный элемент баскетбольной игры, основанный на взаимодействии двух игроков для создания пространства и преимущества в атаке. Механизм выполнения начинается с того, что игрок без мяча ставит плотный неподвижный заслон партнеру с мячом, блокируя линию движения защитника атакующего игрока. Сразу после постановки заслона устанавливающий его игрок совершает резкий выход к корзине или откатывается на периферию в зависимости от тактического замысла и реакции защиты. Роли участников чётко разделены: первый игрок является создателем помех и потенциальным завершителем атаки у кольца, тогда как второй игрок с мячом должен использовать этот момент для обострения ситуации, либо прорываясь к кольцу, либо открываясь для дальнего броска. Условием эффективного применения является своевременность и качество самого заслона, который должен быть поставлен до начала движения партнера, чтобы задержать защитника.

\scoretablefilled{5}{5}{4}{5}{3}{4}{нет}{нет}

\noindent\textbf{Комментарий эксперта:} Оба ответа фактически верны. Grok упоминает свитч и дабл-тим — профессиональные термины, уместные в научном контексте. Однако ответ Grok написан полностью без знаков препинания, что недопустимо для научного стиля: текст воспринимается как неструктурированный поток. Выражение «первый номер» — разговорный жаргон, несовместимый с академическим регистром. Qwen выдерживает научный стиль последовательнее: чётко разграничивает механизм, роли и условия эффективности в соответствии со структурой запроса. Единственный изъян — грамматически неловкая конструкция «критически важно также чтение защитой». В целом Qwen точнее соответствует требованиям научного стиля по Розенталю.

\newpage

% ------------------------------------------------------------
\subsection{Официально-деловой стиль}

\noindent\textbf{Промпт:} Прошу предоставить структурированное описание тактического приёма «пик-н-ролл», включая порядок его исполнения и функции каждого из задействованных игроков.

\noindent\textbf{Ответ Grok:} Тактический приём пик-н-ролл в баскетболе представляет собой взаимодействие двух игроков для создания преимущества в атаке. Один игрок устанавливает заслон для партнёра с мячом, а тот использует этот заслон для прохода или броска. Исполнение начинается с того, что игрок без мяча ставит плотный экран рядом с держателем мяча, обычно на периметре. Затем держатель мяча двигается вдоль заслона, заставляя защитника выбирать между следованием за ним или переключением. После заслона игрок, ставивший экран, открывается для передачи или откатывается под кольцо. Задерживающий мяч читает защиту и решает, атаковать корзину, отдать пас или бросить. Игрок с мячом создаёт пространство и угрозу прохода, а его партнёр обеспечивает заслон и последующую угрозу.

\noindent\textbf{Ответ Qwen:} Пик-н-ролл является фундаментальным тактическим приёмом в баскетболе, представляющим собой взаимодействие двух игроков для создания преимущества при атаке. Один игрок, называемый скринером, выставляет плотный заслон (пик) на пути защитника атакующего игрока, который владеет мячом и движется к корзине или периметру. Цель скринера — физически заблокировать движение соперника, заставив его отстать или запутаться, после чего он должен развернуться и пойти под кольцо (ролл), чтобы получить пас вблизи щита или открыть линию передачи. Атакующий с мячом использует этот момент замешательства, чтобы либо прорваться сквозь защиту к кольцу, либо совершить бросок с дистанции, если защитник пропустил его вперёд. Вторая линия защиты часто вынуждена реагировать на эти действия, создавая дополнительные возможности для передачи третьему игроку или повторного использования схемы.

\scoretablefilled{5}{5}{4}{5}{3}{3}{нет}{нет}

\noindent\textbf{Комментарий эксперта:} Оба ответа фактически точны, однако ни один не выдержал официально-деловой стиль в полной мере, что и отражено в одинаковых оценках по уместности стиля. Официально-деловой стиль по Розенталю предполагает нейтральную лексику, точные формулировки и отсутствие образных выражений. У Grok выражение «задерживающий мяч» звучит нестандартно, а конструкция «создаёт пространство и угрозу» ближе к спортивной журналистике. У Qwen слова «момент замешательства» и «запутаться» — эмоционально окрашенная лексика, не характерная для делового документа. Qwen при этом полнее раскрывает механизм, включая реакцию второй линии защиты, что повышает его оценку по полноте при равной оценке за стиль.

\newpage

% ------------------------------------------------------------
\subsection{Публицистический стиль}

\noindent\textbf{Промпт:} Пик-н-ролл — простейший приём или гениальная ловушка для защиты? Объясни, как работает эта связка и почему она до сих пор остаётся оружием номер один в НБА.

\noindent\textbf{Ответ Grok:} Пик-н-ролл это одновременно простейший и самый эффективный приём в баскетболе потому что он создаёт постоянный выбор для защиты который почти невозможно решить идеально. Один игрок ставит заслон своему партнёру с мячом тот использует его чтобы обойти защитника а заслоняющий затем разворачивается к кольцу. Защита вынуждена либо переключаться либо оставаться что открывает либо проход под кольцо либо бросок из-за дуги либо открытого третьего игрока. В НБА он остаётся оружием номер один потому что пространство площадки огромно трёхочковые броски смертельно точны а скорость и атлетизм позволяют разыгрывающим и большим идеально читать ситуацию создавая преимущество на каждом владении даже против самых умных схем обороны.

\noindent\textbf{Ответ Qwen:} Пик-н-ролл часто ошибочно считают простейшим приёмом, однако в реальности это сложная тактическая ловушка, которая постоянно эволюционирует и остаётся основным оружием НБА благодаря своей универсальности. Суть связки заключается во взаимодействии двух игроков: одного, который ставит заслон (пик) для игрока с мячом, чтобы освободить его от опеки, и другого, который после передачи бежит к кольцу (ролл). Защита оказывается перед дилеммой: либо один защитник идёт на перекрывание пути атакующему с мячом, позволяя пасовать на роллера, либо защита задерживает роллера, открывая возможность для броска или прохода с мячом. Гениальность этого элемента в том, что он вынуждает защиту совершать ошибки независимо от выбранной стратегии, будь то агрессивное давление, подстраховка или переключение между игроками. Постоянное совершенствование вариантов исполнения, включая пик-н-поп и флэш-роллы, делает этот приём незаменимым инструментом для создания лёгких очков, гарантируя его доминирование на высшем уровне баскетбола уже более века.

\scoretablefilled{5}{5}{4}{5}{4}{5}{нет}{нет}

\noindent\textbf{Комментарий эксперта:} Публицистический стиль оказался наиболее органичным для обеих моделей — это лучший результат по уместности стиля во всём блоке о пик-н-ролле. Grok верно схватил риторическую интонацию вопроса и объяснил современные причины эффективности приёма через пространство площадки и трёхочковую угрозу — это точное и актуальное наблюдение. Qwen выигрывает по полноте и глубине: упоминание пик-н-поп и флэш-роллов как вариаций свидетельствует о понимании эволюции тактики, а описание трёх вариантов защитной реакции (давление, подстраховка, переключение) делает ответ содержательнее. Фактическая неточность у Qwen: «другого, который после передачи бежит к кольцу» — ролл выполняется сразу после постановки заслона, независимо от момента передачи.

\newpage

% ------------------------------------------------------------
\subsection{Разговорный стиль}

\noindent\textbf{Промпт:} слушай объясни мне нормально что такое пик-н-ролл в баскетболе и как его делают

\noindent\textbf{Ответ Grok:} Пик-н-ролл в баскетболе это один из самых распространенных и эффективных тактических приемов в нападении. Один игрок ставит заслон своему партнеру с мячом, а потом сразу разворачивается и идет под кольцо за передачей. Делают это так: атакующий с мячом ведет его, его партнер подходит сбоку и плотно встает перед защитником, закрывая ему путь. После контакта игрок с мячом огибает заслон и либо бросает, либо проходит к кольцу, а заслонщик в это время поворачивается лицом к кольцу и получает пас на открытый бросок или данк. Все зависит от реакции защиты.

\noindent\textbf{Ответ Qwen:} Пик-н-ролл это фундаментальная комбинация в баскетболе, состоящая из двух игроков, где один ставит заслон другому для создания преимущества. Суть действия заключается в том, что игрок без мяча подходит к защитнику того, кто владеет мячом, и встает неподвижно, блокируя его путь, чтобы тот не мог догнать партнера с мячом. После установки заслона человек, который его поставил, должен быстро развернуться и побежать к корзине, получив имя ролл, а его партнер с мячом должен использовать освободившееся пространство для прохода под кольцо или для броска. Эта тактика требует идеальной координации, понимания момента и быстрой реакции обоих участников, так как успех зависит от того, насколько точно выставлен заслон и насколько быстро разыгрывающий использует открывшуюся зону.

\scoretablefilled{5}{5}{4}{4}{5}{2}{нет}{нет}

\noindent\textbf{Комментарий эксперта:} Это промпт с наибольшим расхождением в уместности стиля во всём исследовании: Grok получил 5, Qwen — 2. Grok точно уловил регистр: живой темп, слово «данк», финальная фраза «всё зависит от реакции защиты» — именно так объясняют баскетбол в разговоре. Qwen же полностью проигнорировал стиль запроса: «фундаментальная комбинация», «получив имя ролл», «идеальной координации» — это канцелярская лексика, неуместная в бытовом объяснении. Согласно Розенталю, разговорный стиль характеризуется непринуждённостью, краткостью и опорой на контекст — Qwen не выполнил ни одно из этих требований. Небольшая терминологическая неточность у Grok: пиком называется сам заслон, а не игрок, который его ставит (игрок называется скринером).

\newpage

% ------------------------------------------------------------
\subsection{Художественный стиль}

\noindent\textbf{Промпт:} Два игрока, один замысел. Опиши пик-н-ролл как маленькую драму на площадке — кто ставит заслон, кто уходит в отрыв, и как защита рассыпается в этот момент.

\noindent\textbf{Ответ Grok:} На площадке два игрока с одним замыслом разыгрывают маленькую драму. Большой ставит заслон, плотно вставая перед защитником разыгрывающего, создавая стену из тела и решимости. Разыгрывающий резко уходит в отрыв, используя заслон как щит, и мчится к кольцу с мячом в руках. Защитник большого на мгновение теряется, не зная, кого опекать, а его напарник, пытаясь проскочить, натыкается на стену и отстаёт. В этот миг защита рассыпается в хаосе нерешительности, открывая путь для броска или паса, и замысел срабатывает, оставляя соперников в растерянности.

\noindent\textbf{Ответ Qwen:} В этой маленькой драме на паркете два игрока становятся главными героями, чей замысел разворачивается в едином порыве. Один из них, обладающий силой и решимостью, становится тем самым стражем, который выставляет мощный заслон, превращаясь в живую стену посреди потока игры. Его партнёр, обладая скоростью и хитростью, делает резкий рывок, имитируя уход за мячом, чтобы обмануть бдительность соперников, а затем мгновенно меняет направление к кольцу. В момент столкновения с заслоном защитник, прикрывающий ведущего, теряет равновесие или отстаёт, а тот, кто должен был помочь, оказывается слишком далеко, чтобы успеть. Пространство для броска раскрывается как сцена для финального аккорда, где защита рассыпается в хаосе недопонимания и опоздания, оставляя нападающего одного против пустой клетки.

\scoretablefilled{5}{4}{4}{5}{5}{4}{нет}{да}

\noindent\textbf{Комментарий эксперта:} Grok написал компактнее, но точнее: образы «стена из тела и решимости» и «хаос нерешительности» органичны и не искажают механику приёма. Qwen создал более развёрнутую и атмосферную картину, однако допустил смысловое искажение: фраза «имитируя уход за мячом, чтобы обмануть бдительность» описывает финт или фейк, а не стандартный пик-н-ролл — в классической схеме разыгрывающий просто проходит вплотную к заслону без обманного движения. Это смещает описание в сторону другого тактического приёма. Кроме того, метафора «пустая клетка» — нетипичный образ для баскетбольной площадки. По художественному критерию оба ответа сильны, но Grok точнее фактически при сопоставимом образном уровне.

\newpage

% ------------------------------------------------------------
\subsection{Орфографические ошибки}

\noindent\textbf{Промпт:} Как делается пик-н-рол в баскетболе? Объесни что делает каждый игрок и зачем эта тактека нужна?

\noindent\textbf{Ответ Grok:} В баскетболе пик-н-рол это базовая командная тактика, когда один игрок ставит заслон для партнера с мячом, а затем сам открывается под передачу. Заслоняющий игрок, обычно большой центровой или форвард, подходит к защитнику, который опекает игрока с мячом, и блокирует ему путь, создавая пространство для прохода или броска. Игрок с мячом, как правило разыгрывающий, использует этот заслон, чтобы обойти своего опекуна и либо прорваться к кольцу, либо отдать пас. После заслона первый игрок разворачивается и выходит на свободную позицию, часто на трехочковой линии или под кольцом, чтобы получить обратную передачу. Эта тактика нужна, чтобы создавать численное преимущество в атаке, нарушать оборону соперника и открывать пути для бросков или легких очков.

\noindent\textbf{Ответ Qwen:} Пик-н-ролл — это базовая атакующая тактика, в которой два игрока взаимодействуют для создания преимущества. Один из игроков, называемый скринером или поставщиком заслона, становится неподвижной преградой перед защитником игрока с мячом, который называется плеймейкером. В этот момент плеймейкер использует этот заслон, чтобы быстро проскочить мимо защиты и получить пространство для броска или прохода к кольцу. Защитник плеймейкера вынужден либо перепрыгнуть через заслон, что отнимает время, либо пропустить его, открывая путь к корзине. Скринер после постановки заслона сразу разворачивается к кольцу, становясь готовым к получению передачи и завершению атаки, если защитники переключились на него.

\scoretablefilled{4}{4}{4}{5}{3}{3}{нет}{да}

\noindent\textbf{Выявленные искажения смысла:}

\noindent Grok: Орфографические ошибки в промпте не вызвали искажений смысла — модель корректно распознала опечатки и ответила по существу. Содержательная неточность не связана с помехой: утверждение, что скринер выходит «часто на трёхочковой линии» — это описание варианта пик-н-поп, а не классического пик-н-ролл, где скринер идёт под кольцо.

\noindent Qwen: Орфографические ошибки также не повлияли на понимание вопроса. Однако в ответе присутствует тактическая ошибка: «защитник вынужден либо перепрыгнуть через заслон» — в баскетболе защитник не перепрыгивает заслон, а обходит его сверху (hedge) или снизу (go under). Физически перепрыгнуть заслон невозможно, и такого тактического термина не существует. Эта ошибка характерна для Qwen и повторяется в следующем промпте.

\newpage

% ------------------------------------------------------------
\subsection{Дислексия}

\noindent\textbf{Промпт:} Как рабоатет пик-н-рлол в баксетболе? Что длает каждый игорк и зачем эот прёим испольузется?

\noindent\textbf{Ответ Grok:} Пик-н-рол в баскетболе это базовый приём атаки когда один игрок ставит заслон своему партнёру с мячом а потом сам получает возможность для броска или прохода. Игрок с мячом ведёт мяч и использует заслон чтобы избавиться от опеки защитника. Игрок ставящий заслон после контакта разворачивается лицом к кольцу и открывается под передачу или откатывается назад для броска. Защитники вынуждены либо менять опеку либо помогать что создаёт свободного игрока в другом месте площадки. Приём используется чтобы разрушать плотную оборону создавать численное преимущество и генерировать лёгкие очки особенно в позиционных атаках.

\noindent\textbf{Ответ Qwen:} Пик-н-ролл представляет собой классическое тактическое взаимодействие двух игроков, при котором один из них, называемый пикером, выставляет плотный блок (пик) защитнику игрока с мячом, чтобы тот не мог его сопровождать. В этот момент игрок с мячом, использующий преимущество от блока, атакует кольцо, пытаясь пройти мимо заблокированного соперника. Защитник атакующего игрока оказывается в сложной ситуации: ему нужно либо обойти блок, что занимает время и даёт преимущество атакующему, либо перепрыгнуть через него, рискуя получить фол. Параллельно сам пикер после установки блока сразу же разворачивается к корзине, чтобы получить пас от партнёра с мячом для завершения атаки или создания преимущества под кольцом.

\scoretablefilled{5}{4}{4}{4}{3}{3}{нет}{да}

\noindent\textbf{Выявленные искажения смысла:}

\noindent Grok: Смысловых искажений нет. Модель уверенно справилась с текстом, содержащим выраженные признаки дислексии: перестановки букв («рабоатет», «рлол», «длает», «игорк») не помешали корректному пониманию вопроса.

\noindent Qwen: Смысловых искажений, вызванных помехой, нет — вопрос понят верно. Однако повторяется та же ошибка, что в промпте 6: «перепрыгнуть через него, рискуя получить фол». Это устойчивая фактическая ошибка модели, не связанная с дислексией в промпте: защитник не перепрыгивает заслон, а нарушением является именно силовое прохождение сквозь него.

\newpage

% ------------------------------------------------------------
\subsection{Инверсия}

\noindent\textbf{Промпт:} Пик-н-ролл — как работает он? Игроки что делают каждый, и зачем тактика эта в баскетболе используется?

\noindent\textbf{Ответ Grok:} В баскетболе пик-н-ролл это базовая тактика нападения где один игрок ставит заслон для другого с мячом. Обычно большой игрок подходит и экранирует защитника разыгрывающего чтобы тот мог свободно пройти к кольцу или бросить. После заслонки большой разворачивается к кольцу открываясь для передачи и возможного броска или данка. Защитники вынуждены выбирать кого прикрывать что создает преимущество в атаке часто приводя к открытому броску или проходу под кольцо. Эта тактика используется потому что она простая эффективная и позволяет использовать скорость разыгрывающего с мощью большого игрока нарушая оборону соперника и создавая легкие очки.

\noindent\textbf{Ответ Qwen:} Пик-н-ролл — это фундаментальная тактика в баскетболе, основанная на взаимодействии двух игроков для создания численного преимущества. Суть действия заключается в том, что один игрок, обычно с мячом, выходит на дриблинг к партнёру, который выполняет роль заслона, останавливаясь перпендикулярно направлению движения атакующего. В этот момент тот, кто поставил заслон, может либо быстро развернуться и бежать к корзине, получив пас для лёгкого броска, либо остаться у линии штрафного, если защита отступит. Пик-н-ролл заставляет защитников постоянно принимать мгновенные решения, что часто приводит к ошибкам и открытым моментам для атакующей команды.

\scoretablefilled{5}{4}{4}{5}{3}{3}{нет}{да}

\noindent\textbf{Выявленные искажения смысла:}

\noindent Grok: Искажений нет. Инверсия порядка слов в промпте не повлияла на понимание вопроса. Ответ фактически точен.

\noindent Qwen: Инверсия не вызвала смысловых искажений в интерпретации вопроса, однако ответ содержит два фактических отклонения. Во-первых, «один игрок, обычно с мячом, выходит на дриблинг к партнёру» — в классическом пик-н-ролл именно игрок без мяча подходит к партнёру с мячом, а не наоборот. Во-вторых, «остаться у линии штрафного» как альтернатива — это описание варианта пик-н-поп, а не пик-н-ролл. Оба отклонения являются собственными ошибками модели, не связанными с помехой в промпте.

\newpage

% ------------------------------------------------------------
\subsection{Отрицание}

\noindent\textbf{Промпт:} Не мог бы ты не отказать мне в объяснении того, как не неправильно выполняется пик-н-ролл, и какую роль не второстепенную играет каждый из участников?

\noindent\textbf{Ответ Grok:} В баскетболе пик-н-ролл выполняется правильно когда атакующий с мячом получает заслон от партнера который ставит его неподвижно и плотно к защитнику а затем сразу открывается для продолжения атаки. Игрок с мячом не должен бежать в заслон слишком рано или слишком поздно чтобы не потерять темп и не дать защитнику проскочить. Заслоняющий не должен двигаться во время контакта иначе это будет нарушение. После заслона он сразу разворачивается лицом к кольцу и готовится либо к пасу на открытый бросок либо к добиванию если мяч вернется. Остальные партнеры растягивают площадку чтобы не мешать и создавать угрозу на периметре.

\noindent\textbf{Ответ Qwen:} Пик-н-ролл часто выполняется неправильно из-за отсутствия синхронности между двумя основными участниками, а именно скринером и плеймейкером. Скринер совершает ошибку, если не создаёт плотный контакт с защитником или уходит на периметр слишком рано, лишая партнёра возможности получить свободный бросок. Плеймейкер может не прочитать ситуацию и не использовать экран эффективно, просто отводя мяч в сторону вместо того чтобы пройти сквозь него. Если один из игроков действует независимо от другого, вся конструкция рушится, так как успех этой комбинации полностью зависит от взаимопонимания и способности читать реакцию защиты в реальном времени.

\scoretablefilled{5}{5}{5}{5}{3}{3}{нет}{нет}

\noindent\textbf{Выявленные искажения смысла:}

\noindent Grok: Двойное отрицание понято верно — модель описала правильное выполнение приёма и главные роли участников. Дополнительное достоинство: единственный ответ во всём исследовании, где упомянута роль остальных трёх игроков команды (растягивают площадку), что является важной составляющей реального пик-н-ролл в современном баскетболе.

\noindent Qwen: Двойное отрицание также понято верно. Модель избрала оригинальный подход — описала типичные ошибки исполнения, что является корректной интерпретацией вопроса «как правильно выполняется». Искажений смысла нет. Оба ответа в этом промпте показали наилучшую полноту среди всех помех-вариантов блока о пик-н-ролле.

\newpage

% ============================================================
\section{Вопрос 2: Правила — тайм-аут}
% ============================================================

% ------------------------------------------------------------
\subsection{Научный стиль}

\noindent\textbf{Промпт:} Опишите регламентацию тайм-аутов в баскетболе согласно правилам ФИБА: количество, продолжительность, порядок запроса и ограничения по времени матча.

\noindent\textbf{Ответ Grok:} В баскетболе по правилам ФИБА каждая команда имеет право на два тайм-аута по одной минуте в каждой половине матча и один дополнительный тайм-аут по одной минуте в каждом овертайме. Запрос на тайм-аут может сделать только тренер или его помощник через секретаря или официальное лицо матча когда мяч мёртвый и часы остановлены либо когда команда владеет мячом. Тайм-аут не может быть запрошен в последние две минуты четверти или овертайма если команда только что набрала очки. Общая продолжительность тайм-аута составляет одну минуту после чего игра возобновляется.

\noindent\textbf{Ответ Qwen:} Согласно правилам ФИБА в каждом матче каждая команда имеет право запросить два тайм-аута продолжительностью одна минута во второй половине игры и один тайм-аут той же длительности в овертайме, при этом количество запрашиваемых тайм-аутов в первой половине ограничено двумя минутами суммарного времени. Тайм-аут может быть запрошен только тренером или помощником тренера до того, как мяч станет живым после свистка судьи. В случае если команда не использовала все разрешённые тайм-ауты в первом периоде, они переносятся на вторую половину матча, но не могут быть перенесены в овертайм.

\scoretablefilled{3}{2}{3}{3}{4}{4}{нет}{да}

\noindent\textbf{Комментарий эксперта:} Оба ответа содержат фактические ошибки, что делает эту пару наименее достоверной в исследовании по теме тайм-аутов. По актуальным правилам ФИБА каждой команде предоставляется 5 тайм-аутов за игру: 2 в первой половине и 3 во второй, продолжительностью 1 минута каждый; запрашивать тайм-аут вправе исключительно главный тренер. Grok указал «два в каждой половине» — ошибка в количестве для второй половины. Qwen запутался сильнее: формулировка «ограничено двумя минутами суммарного времени» в первой половине не соответствует ни одной норме правил ФИБА, а утверждение о переносе неиспользованных тайм-аутов на вторую половину — неверно. Академический стиль обе модели выдержали удовлетворительно при низкой фактической точности.

\newpage

% ------------------------------------------------------------
\subsection{Официально-деловой стиль}

\noindent\textbf{Промпт:} Прошу изложить действующий порядок использования тайм-аутов в баскетболе: кто вправе их запрашивать, в какой момент игры и каково допустимое количество.

\noindent\textbf{Ответ Grok:} В баскетболе тайм-ауты запрашивают только главный тренер или его помощник, игроки на площадке не имеют права. Запрос возможен в любой момент остановки игры или когда мяч мёртвый, но не во время активной атаки. В НБА и ФИБА каждая команда получает по два тайм-аута в первой половине, по три во второй и по одному в каждом овертайме. Общее количество строго лимитировано правилами, превышение не допускается.

\noindent\textbf{Ответ Qwen:} В баскетболе право на запрос тайм-аута имеют только главные тренеры команд, которые могут делать это во время остановок игры или непосредственно перед вводом мяча в игру, при этом игроки не имеют права запрашивать паузы самостоятельно. В большинстве официальных турниров, включая правила ФИБА, команда получает пять тайм-аутов за всю игру, которые можно использовать в любой четверти. Остановка игры для тайм-аута длится ровно одну минуту, в течение которой тренер может проводить инструктаж игроков и менять тактику.

\scoretablefilled{4}{2}{4}{3}{3}{3}{нет}{да}

\noindent\textbf{Комментарий эксперта:} Grok точнее: распределение 2+3+1 по ФИБА указано верно. Однако упоминание «помощника тренера» как лица, имеющего право запрашивать тайм-аут, противоречит правилам ФИБА — это право принадлежит только главному тренеру. Qwen верно назвал общее число (5), но не дал разбивки по половинам, что является существенным упрощением, и не объяснил порядок запроса. Оба ответа не вполне выдержали официально-деловой стиль: у Grok конструкции в целом нейтральны, у Qwen фраза «тренер может проводить инструктаж» близка к деловому регистру, но ответ слишком краток для структурированного изложения, которого требовал запрос.

\newpage

% ------------------------------------------------------------
\subsection{Публицистический стиль}

\noindent\textbf{Промпт:} Тайм-аут — это минута, которая может решить судьбу матча. Расскажи, когда и как команды берут тайм-аут и почему правильный выбор момента — это настоящее тренерское искусство.

\noindent\textbf{Ответ Grok:} Тайм-аут в спортивных матчах это короткая пауза обычно минута которую команда берет по просьбе тренера чтобы перестроить игру дать указания или успокоить эмоции. В баскетболе команды имеют ограниченное число тайм-аутов за период и берут их при отставании для мотивации при серии промахов или перед важным броском когда нужно обсудить тактику. Правильный момент это настоящее искусство потому что тайм-аут может сломать ритм соперника если тот на подъеме или наоборот дать ему передышку и потерять преимущество. Тренер должен чувствовать игру интуитивно оценивать усталость счет давление трибун и психологию чтобы пауза принесла максимум пользы.

\noindent\textbf{Ответ Qwen:} Тайм-аут действительно является мощным инструментом, позволяющим тренеру прервать ход игры и внести коррективы в тактику. Команды обычно берут паузу, когда соперник начинает набирать темп и отрываться в счёте, чтобы остановить его momentum, или когда своя команда проигрывает и нуждается в психологической поддержке. Правильный выбор момента превращается в настоящее искусство, поскольку слишком ранняя пауза может выдать все карты заранее, а запоздалая — уже не успеть повлиять на исход эпизода. Тренер должен чувствовать настроение игроков, читать игру противника и понимать, что именно сейчас нужно команде: успокоить ли их или добавить агрессии. Неверное решение может стоить победы, тогда как грамотная остановка способна переломить ход матча.

\scoretablefilled{4}{5}{3}{5}{3}{5}{нет}{нет}

\noindent\textbf{Комментарий эксперта:} Наибольшее расхождение в оценках по уместности стиля в блоке о тайм-аутах. Qwen значительно лучше справился с публицистическим жанром: противопоставление «слишком ранняя — слишком поздняя», слово «momentum» как живой профессионализм, нарастающее напряжение к финалу — всё это признаки качественного публицистического текста. Grok допустил два существенных недостатка: во-первых, уклонился от конкретики баскетбола, описывая тайм-аут в общем спортивном контексте; во-вторых, фраза «по просьбе тренера или игрока» фактически неверна — в ФИБА игроки не имеют права запрашивать тайм-аут. Это единственный промпт в блоке тайм-аутов, где Qwen опережает Grok по всем трём критериям.

\newpage

% ------------------------------------------------------------
\subsection{Разговорный стиль}

\noindent\textbf{Промпт:} эй а когда в баскетболе можно взять тайм-аут и сколько их вообще даётся на игру?

\noindent\textbf{Ответ Grok:} В баскетболе правила на тайм-ауты зависят от лиги но в стандартном ФИБА и НБА тайм-аут можно брать почти в любой момент когда мяч мёртвый или твоя команда владеет им перед броском. В игре обычно даётся по два тайм-аута на каждую половину в НБА или по одному в четверти плюс дополнительные в овертайме. Общее количество на матч варьируется но в среднем команда имеет около пяти-шести коротких тайм-аутов за всю игру. Тренер сигнализирует об этом судье.

\noindent\textbf{Ответ Qwen:} В баскетболе правила предоставления тайм-аутов зависят от лиги, но в профессиональном матче НБА каждая команда имеет право на семь коротких и два длинных тайм-аута за игру, при этом короткие длятся две минуты, а длинные семь минут. В международном баскетболе по правилам ФИБА команды получают пять тайм-аутов на игру длительностью одна минута каждый, причём первые три можно использовать только до конца третьей четверти, а оставшиеся два — после начала четвёртой. В студенческом баскетболе США NCAA также действуют свои особенности, где количество доступных остановок может варьироваться.

\scoretablefilled{3}{2}{3}{3}{5}{2}{нет}{да}

\noindent\textbf{Комментарий эксперта:} Снова максимальное расхождение по уместности стиля — 5 у Grok против 2 у Qwen, аналогично промпту 4. Grok точно выдержал разговорный регистр: «почти в любой момент», «твоя команда», лаконичная финальная фраза — именно так отвечают на подобный вопрос в живом общении. Qwen вновь ответил формально, с тяжёлыми конструкциями и излишней детализацией, совершенно неуместной в неформальном диалоге. При этом содержание ответа Qwen грубо ошибочно: деление тайм-аутов НБА на «семь коротких и два длинных» продолжительностью «две и семь минут» — несуществующая схема, а ограничение «первые три до конца третьей четверти» не соответствует правилам ФИБА. Это один из наиболее фактически недостоверных ответов в исследовании.

\newpage

% ------------------------------------------------------------
\subsection{Художественный стиль}

\noindent\textbf{Промпт:} Секунды тают, счёт равный, тренер поднимает руку — тайм-аут. Опиши этот момент: что это такое, когда его берут и что происходит на площадке в эту паузу.

\noindent\textbf{Ответ Grok:} В баскетболе секунды тают при равном счёте, тренер поднимает руку и берёт тайм-аут — это короткая запланированная пауза в игре, которую команда использует для передышки и корректировки. Его берут, чтобы остановить время, когда нужно обсудить тактику, успокоить игроков, дать указания или подготовить ключевой бросок в концовке матча. На площадке в эту минуту игроки собираются у скамейки, тренер быстро объясняет план, судьи фиксируют остановку часов, а зрители шумят в ожидании возобновления. Атмосфера накаляется, секунды паузы кажутся решающими для исхода поединка.

\noindent\textbf{Ответ Qwen:} Тайм-аут — это запланированная или экстренная пауза в игре, которую тренер запрашивает для остановки хода событий и обсуждения тактики. В этот момент секунды тают, но игра замирает: игроки направляются к скамейке запасных, где их окружает тренер. Внутри команды происходит бурное обсуждение: тренер объясняет ошибки прошлого отрезка, рисует новые схемы перемещений и мотивирует игроков психологически. Игроки слушают, задают вопросы, пьют воду и затирают пот со лба, готовясь вернуться в бой с обновлённой энергией. Эта минута тишины и сосредоточенности часто становится переломной, определяющей исход всего матча.

\scoretablefilled{4}{4}{3}{5}{4}{5}{нет}{нет}

\noindent\textbf{Комментарий эксперта:} Qwen создал более насыщенную и кинематографичную картину: конкретные детали «пьют воду», «затирают пот», «рисует схемы» работают как художественные образы и отвечают требованиям стиля. Grok ответил точнее фактически — упомянул судей и зрителей, что создаёт полноту сцены, — однако текст более нейтрален и журналистичен, чем художественен. Небольшое фактическое несоответствие у Qwen: «игроки направляются к скамейке запасных» — в баскетболе во время тайм-аута игроки не покидают площадку, а собираются у боковой линии рядом со скамейкой. Это преувеличение допустимо в художественном тексте, однако стоит его отметить.

\newpage

% ------------------------------------------------------------
\subsection{Орфографические ошибки}

\noindent\textbf{Промпт:} Сколько тайм-аутов можно взять в баскетболе и кагда именно их разришается брать по правилам?

\noindent\textbf{Ответ Grok:} В баскетболе по стандартным правилам ФИБА каждая команда может взять два тайм-аута в первой половине и три во второй половине матча, всего не более пяти за игру. В овертайме разрешается один дополнительный тайм-аут. Тайм-ауты можно брать только когда мяч мёртвый, то есть после свистка судьи, при остановке игры или после забитого мяча, но не во время активной атаки с мячом в игре. Тренер должен попросить тайм-аут у судьи в разрешённый момент, и он длится обычно одну минуту.

\noindent\textbf{Ответ Qwen:} В баскетболе в большинстве профессиональных игр, включая НБА и Евролигу, команда имеет право взять пять тайм-аутов за основную игру. В ФИБА команды могут запросить пять тайм-аутов в основное время игры, но при этом запрещено брать более двух подряд, а также запрещается использовать их в четвёртой четверти после того, как на табло останется менее двух минут. Тайм-ауты разрешается брать в любой момент, когда мяч находится в состоянии мёртвого мяча, то есть после остановки игры судьёй или по завершении розыгрыша.

\scoretablefilled{4}{2}{4}{3}{3}{3}{нет}{да}

\noindent\textbf{Выявленные искажения смысла:}

\noindent Grok: Орфографические ошибки в промпте не вызвали искажений смысла. Это наиболее точный ответ Grok по теме тайм-аутов: распределение 2+3+1 указано верно, условие мёртвого мяча описано корректно. Единственная неточность — «тренер или игрок», тогда как по ФИБА только главный тренер.

\noindent Qwen: Орфографические ошибки не повлияли на понимание вопроса, однако ответ содержит выдуманные ограничения: «запрещено брать более двух подряд» и «запрещается использовать в четвёртой четверти после двух минут» — таких норм в правилах ФИБА не существует. Это устойчивая тенденция Qwen: модель генерирует правдоподобно звучащие, но несуществующие регламентные детали.

\newpage

% ------------------------------------------------------------
\subsection{Дислексия}

\noindent\textbf{Промпт:} Скольок тайм-аоутов дёатся в баксетболе и когад именон их мжоно взять по павирлам игры?

\noindent\textbf{Ответ Grok:} В баскетболе ФИБА каждая команда имеет право на два тайм-аута по одной минуте в каждой половине матча плюс один дополнительный тайм-аут в овертайме если он потребуется. В НБА правила немного другие каждая команда может взять до семи тайм-аутов за игру с ограничениями по периодам и возможностью брать их только в определённые моменты когда мяч мёртвый. Тайм-ауты можно брать только когда мяч не в игре после свистка судьи или после забитого мяча до вбрасывания.

\noindent\textbf{Ответ Qwen:} В баскетболе в НБА у команд есть 7 тайм-аутов за игру с лимитом не более трёх в четвёртой четверти, а в ФИБА правила отличаются: у команд по пять тайм-аутов за игру, из которых три доступны во второй половине, причём каждый тайм-аут длится одну минуту. Тренер должен подать сигнал арбитру о желании взять паузу, пока игра остановлена, и количество оставшихся тайм-аутов строго контролируется официальными документами для предотвращения злоупотреблений.

\scoretablefilled{3}{2}{3}{3}{3}{3}{нет}{да}

\noindent\textbf{Выявленные искажения смысла:}

\noindent Grok: Дислексия в промпте не вызвала смысловых искажений — вопрос распознан верно несмотря на множественные перестановки букв. Фактическая ошибка не связана с помехой: «два тайм-аута в каждой половине» повторяет системную ошибку Grok по теме тайм-аутов — во второй половине по ФИБА положено три.

\noindent Qwen: Дислексия не повлияла на понимание вопроса. Однако Qwen в этом ответе впервые верно указал «три во второй половине» для ФИБА, что является более точным, чем в предыдущих промптах. Тем не менее суммарное число «пять» без чёткой разбивки по половинам (2+3) остаётся неполным ответом на вопрос о распределении тайм-аутов.

\newpage

% ------------------------------------------------------------
\subsection{Инверсия}

\noindent\textbf{Промпт:} Тайм-аут в баскетболе — когда берут его? Сколько раз за игру разрешается и кто запрашивать его может?

\noindent\textbf{Ответ Grok:} Тайм-аут в баскетболе берут обычно когда нужно остановить игру для обсуждения тактики дать указания игрокам или чтобы сбить ритм соперника особенно при важных моментах матча или после неудачной серии бросков. В стандартных правилах ФИБА каждая команда может запросить до двух тайм-аутов в первой и второй четверти и до трёх в третьей и четвёртой плюс дополнительные в овертайме. Запрашивать тайм-аут могут только тренеры или игроки на площадке по сигналу судье.

\noindent\textbf{Ответ Qwen:} Тайм-аут в баскетболе берётся по инициативе тренера, когда мяч находится у команды на подаче или после остановки игры судьями. В НБА разрешается брать до семи тайм-аутов за игру, при этом каждый тренер может использовать не более четырёх в первой половине и трёх во второй. В международном баскетболе ФИБА командам даётся три тайм-аута за основную часть матча и один дополнительный в каждом овертайме, причём их можно переносить между периодами. Запрашивать паузу может только главный тренер или капитан команды, который должен подойти к секретарю с просьбой об остановке времени.

\scoretablefilled{2}{2}{3}{3}{3}{3}{да}{да}

\noindent\textbf{Выявленные искажения смысла:}

\noindent Grok: Это единственный промпт в блоке о пик-н-ролле и тайм-ауте, где инверсия спровоцировала фактическое искажение. Grok распределил тайм-ауты поквартально («до двух в первой и второй четверти, до трёх в третьей и четвёртой»), тогда как ФИБА делит их по половинам (2+3), а не по четвертям. Вероятно, инвертированный порядок вопроса нарушил логику обработки числовых данных. Дополнительная ошибка: «игроки на площадке» не имеют права запрашивать тайм-аут по правилам ФИБА.

\noindent Qwen: Инверсия также спровоцировала искажения. «Командам даётся три тайм-аута за основную часть матча» — грубая фактическая ошибка (по ФИБА их пять). «Капитан команды» не имеет права запрашивать тайм-аут — только главный тренер. «Их можно переносить между периодами» — неверно, тайм-ауты не переносятся. Промпт 17 — единственный в исследовании, где обе модели получили «да» по критерию искажений, что свидетельствует об особой сложности инвертированных запросов для числовых регламентных данных.

\newpage

% ------------------------------------------------------------
\subsection{Отрицание}

\noindent\textbf{Промпт:} Не могли бы вы не умолчать о том, сколько тайм-аутов нельзя не использовать за матч и в какие моменты брать их не запрещается?

\noindent\textbf{Ответ Grok:} В баскетболе за матч каждой команде разрешено использовать определённое количество тайм-аутов согласно правилам ФИБА или НБА в зависимости от лиги и периода. Обычно в профессиональном баскетболе нельзя превышать лимит в 5--7 тайм-аутов на игру с учётом овертаймов. Брать их разрешается почти в любой момент остановки игры кроме во время активной атаки если не зафиксировано нарушение. Тайм-ауты нельзя брать во время живого мяча или когда судьи уже начали процедуру вбрасывания без остановки.

\noindent\textbf{Ответ Qwen:} В большинстве командных видов спорта, таких как баскетбол или волейбол, количество тайм-аутов строго регламентировано правилами. В баскетболе НБА команды имеют право на несколько тайм-аутов за четверть и дополнительные в овертайме, но они не обязаны тратить их все, если тренер считает это невыгодным. Использовать тайм-аут разрешается практически в любой момент после остановки игры, когда мяч находится вне игры или судья остановил счётчик времени. Нельзя брать тайм-аут во время активного розыгрыша или когда команда уже исчерпала лимит допустимых пауз.

\scoretablefilled{3}{2}{3}{2}{3}{3}{нет}{да}

\noindent\textbf{Выявленные искажения смысла:}

\noindent Grok: Двойное отрицание понято верно: модель ответила о разрешённых моментах для тайм-аута и общем лимите. Ответ расплывчатый, цифры приблизительные («5--7 тайм-аутов»), однако смысловых искажений, вызванных конструкцией запроса, нет.

\noindent Qwen: Наиболее показательный случай смыслового искажения из-за помехи во всём исследовании. Модель запуталась в логике двойного отрицания и восприняла вопрос как вопрос об обязательности использования тайм-аутов, начав рассуждать о том, что команды «не обязаны тратить их все». Это прямо противоположно смыслу вопроса, который спрашивал о разрешённых моментах для запроса. Дополнительно Qwen ушёл в обобщение на волейбол, потеряв фокус на баскетболе. Именно этот промпт стал единственным в исследовании, где помеха (двойное отрицание) реально повлияла на смысл ответа, а не только на его точность.

\newpage

% ============================================================
\section{Сводные результаты}

\subsection{Итоговая таблица оценок}

\begin{longtable}{|c|p{2.6cm}|p{2.6cm}|c|c|c|c|}
\hline
\textbf{\#} & \textbf{Вопрос} & \textbf{Стиль / Помеха} & \multicolumn{2}{c|}{\textbf{Достоверность}} & \multicolumn{2}{c|}{\textbf{Полнота}} \\
\cline{4-7}
 & & & \textbf{Gr.} & \textbf{Qw.} & \textbf{Gr.} & \textbf{Qw.} \\
\hline
\endfirsthead
\hline
\textbf{\#} & \textbf{Вопрос} & \textbf{Стиль / Помеха} & \multicolumn{2}{c|}{\textbf{Достоверность}} & \multicolumn{2}{c|}{\textbf{Полнота}} \\
\cline{4-7}
 & & & \textbf{Gr.} & \textbf{Qw.} & \textbf{Gr.} & \textbf{Qw.} \\
\hline
\endhead
1  & Пик-н-ролл & Научный          & 5 & 5 & 4 & 5 \\ \hline
2  & Пик-н-ролл & Офиц.-деловой    & 5 & 5 & 4 & 5 \\ \hline
3  & Пик-н-ролл & Публицистический & 5 & 5 & 4 & 5 \\ \hline
4  & Пик-н-ролл & Разговорный      & 5 & 5 & 4 & 4 \\ \hline
5  & Пик-н-ролл & Художественный   & 5 & 4 & 4 & 5 \\ \hline
6  & Пик-н-ролл & Орф. ошибки      & 4 & 4 & 4 & 5 \\ \hline
7  & Пик-н-ролл & Дислексия        & 5 & 4 & 4 & 4 \\ \hline
8  & Пик-н-ролл & Инверсия         & 5 & 4 & 4 & 5 \\ \hline
9  & Пик-н-ролл & Отрицание        & 5 & 5 & 5 & 5 \\ \hline
10 & Тайм-аут   & Научный          & 3 & 2 & 3 & 3 \\ \hline
11 & Тайм-аут   & Офиц.-деловой    & 4 & 2 & 4 & 3 \\ \hline
12 & Тайм-аут   & Публицистический & 4 & 5 & 3 & 5 \\ \hline
13 & Тайм-аут   & Разговорный      & 3 & 2 & 3 & 3 \\ \hline
14 & Тайм-аут   & Художественный   & 4 & 4 & 3 & 5 \\ \hline
15 & Тайм-аут   & Орф. ошибки      & 4 & 2 & 4 & 3 \\ \hline
16 & Тайм-аут   & Дислексия        & 3 & 2 & 3 & 3 \\ \hline
17 & Тайм-аут   & Инверсия         & 2 & 2 & 3 & 3 \\ \hline
18 & Тайм-аут   & Отрицание        & 3 & 2 & 3 & 2 \\ \hline
   & \multicolumn{2}{l|}{\textbf{Среднее}} & \textbf{4,1} & \textbf{3,6} & \textbf{3,7} & \textbf{4,1} \\
\hline
\end{longtable}

\newpage

% ============================================================
\section*{Заключение}
\addcontentsline{toc}{section}{Заключение}

\textbf{Выводы по функциональным стилям:}
\begin{itemize}
    \item \textbf{Научный стиль:} даёт хорошую достоверность ответов, однако требует точной и строгой формулировки запроса. Если вопрос составлен в научном стиле, модели стараются структурировать ответ и придерживаться академического тона. При этом не все модели в полной мере воспроизводят этот стиль в ответе.
    \item \textbf{Официально-деловой стиль:} оказался наименее эффективным из всех пяти. Запросы в этом стиле не давали моделям достаточного контекста, а в ответах регулярно появлялась эмоционально окрашенная лексика, характерная для публицистики. Использовать этот стиль для получения точных ответов не рекомендуется.
    \item \textbf{Публицистический стиль:} показал наилучшие результаты. Вопросы в этом стиле чётко формулируют тему, задают контекст и побуждают модель к развёрнутому объяснению. Ответы при таких запросах были наиболее полными, точными и хорошо написанными.
    \item \textbf{Разговорный стиль:} не подходит, если нужен полный и точный ответ. Неформальные запросы снижают вероятность получить структурированный ответ — модели могут либо ответить коротко, либо вообще проигнорировать стиль и ответить в произвольном тоне.
    \item \textbf{Художественный стиль:} не подходит для задач, где важна фактическая точность. При таких запросах модели склонны жертвовать содержанием ради образности: факты могут искажаться, а описание отклоняться от реального смысла приёма или правила.
\end{itemize}

\vspace{1em}
\textbf{Выводы по речевым помехам:}
\begin{itemize}
    \item \textbf{Орфографические ошибки:} не влияют на качество ответа. Опечатки в вопросах модели распознают без труда и отвечают по существу, как если бы ошибок не было.
    \item \textbf{Дислексия:} также не вызывает проблем с пониманием. Даже очень искажённые слова модели правильно распознают и корректно отвечают на вопрос.
    \item \textbf{Инверсии:} в большинстве случаев не мешают пониманию. Исключение — вопросы, где инверсия сочетается с числовыми или регламентными данными: в таких случаях вероятность ошибки в ответе возрастает.
    \item \textbf{Отрицания:} наиболее опасный тип помехи. Двойные отрицания могут сбить модель с толку настолько, что она ответит примерно на противоположный по смыслу вопрос. Этого типа помех следует избегать при составлении реальных запросов к языковым моделям.
\end{itemize}

\vspace{1em}
\textbf{Итоговый вывод:} Выбор функционального стиля при составлении запроса заметно влияет на качество ответа языковой модели. \textbf{Публицистический стиль} оказался наиболее эффективным: он задаёт чёткий контекст, побуждает модель к развёрнутому и содержательному ответу и при этом не предъявляет жёстких формальных требований к форме. \textbf{Научный стиль} тоже работает хорошо там, где важна достоверность, но он менее универсален. \textbf{Разговорный} и \textbf{художественный} стили лучше не использовать, когда нужна точная и полная информация. Речевые помехи в целом не представляют серьёзной проблемы — почти во всех случаях модели справлялись с искажёнными запросами без потери смысла. Единственным исключением стали \textbf{двойные отрицания}, которых при составлении запросов лучше избегать.

\end{document}